\chapter{Introduction}

\section{What an AMR Is}

An Autonomous Mobile Robot (AMR) is a robot that navigates an environment
without fixed infrastructure -- no magnetic strips and no rails -- using
sensors such as LiDAR and cameras together with SLAM (Simultaneous
Localization and Mapping) algorithms to move safely around obstacles and
people. This is what separates an AMR from the older Automated Guided Vehicle
(AGV), which can only follow pre-defined physical paths. Over the past decade
AMRs have moved out of research labs into commercial use across manufacturing,
warehousing, healthcare, hospitality and security \cite{standardbots_amr,valuates_agv_amr}.

\section{Why the Market Is Growing}

Three forces are driving global adoption:

\begin{itemize}
    \item \textbf{Labor shortages and rising labor costs} in mature economies
          (the US, Europe, China's coastal regions), which make automation of
          repetitive transport and inspection tasks economically attractive.
    \item \textbf{Falling sensor and compute costs} (LiDAR, depth cameras,
          embedded GPUs), which have made AMR hardware progressively cheaper to
          build over the last five years.
    \item \textbf{Maturing open-source software} (ROS / ROS2 and the Nav2
          navigation stack), which has lowered the barrier to building a working
          navigation system without each company writing one from scratch.
\end{itemize}

A second, more recent trend is \textbf{modularity}: instead of selling a
single-purpose robot, manufacturers increasingly sell a base platform that
carries different payloads -- a delivery box, a robotic arm, a sensor pack --
depending on the client's task \cite{fdatabot_amr2026,robotnik_kairos}. This is
the segment this document focuses on, because it is the direction CyberMech
Systems' own platform is being designed in: one mobile base accepting
interchangeable modules.

\section{Scope and Method}
 
This study does \emph{not} assume that what works internationally will
automatically work in Tunisia. A large part of the analysis below is dedicated
specifically to testing that assumption against the Tunisian economic context --
labor costs, company size, import dependency and procurement habits -- rather
than treating global AMR growth figures as directly transferable. Concretely,
the report covers the overall market and a sector-by-sector landscape
(Chapters~2--3), the competitor field (Chapter~4), target clients, pricing,
pain points and use cases (Chapters~5--8), and a consolidated set of
opportunities and a final recommendation (Chapters~9--10).

\textit{Note on method.} This chapter and those that follow consolidate two
market-research efforts that the team produced in parallel: a global/regional
sizing study built around a four-module platform framework, and a
Tunisia-focused study built around local unit economics. Where the two agreed,
the finding is reported plainly; where they differed, the stronger-evidenced
position is taken and the disagreement is noted.
